\subsection{Prediction 7: Why Everyone Loses When Everyone Optimizes}
\label{subsec:pred7-multilevel}

%%%%%%%%%%%%%%%%%%%%%%%%%%%%%%%%%%%%%%%%%%%%%%%%%%%%%%%%%%%%%%%%%%%%%%%%%%%%%%%
% CHAPTER 11.7: Narrative Treatment of Prediction 7
% Multi-Level Coherence Conflicts
%%%%%%%%%%%%%%%%%%%%%%%%%%%%%%%%%%%%%%%%%%%%%%%%%%%%%%%%%%%%%%%%%%%%%%%%%%%%%%%

\subsubsection{The Climate Summit Paradox}

December 2015. Paris. World leaders gathered for COP21.

Every nation in the room wanted to solve climate change. Polls showed \
overwhelming public support for action. The science was clear. The \
economics, increasingly favorable. Renewable energy costs were plummeting. \
The moment seemed right.

And yet.

The agreement that emerged---celebrated as historic---committed nations \
to goals they systematically miss. Global emissions have risen, not fallen, \
since Paris. Temperature targets that seemed ambitious are now nearly \
impossible. The window for limiting warming to 1.5°C has all but closed.

What went wrong?

Not corruption. Not ignorance. Not even short-termism, exactly.

What went wrong is that everyone optimized---and optimization at every \
level produces outcomes optimal at none.

\subsubsection{The Three-Level Problem}

Consider the structure of climate negotiations:

\paragraph{Level 1: Individuals.}

Maria drives a fuel-efficient car, recycles, and votes for green candidates. \
She feels good about her choices. Her personal coherence---her sense that \
her actions align with her values---is high.

But Maria also flies to visit family, heats her home in winter, and buys \
products shipped from China. When asked about these choices, she has good \
reasons: family is important, heating isn't optional, and alternatives \
aren't available.

Maria optimizes her own coherence, $K_{ind}$. She feels like she's doing \
her part.

\paragraph{Level 2: Nations.}

Germany invested heavily in solar power and announced ambitious targets. \
German politicians feel good about their leadership. Germany's national \
coherence---its sense of being a responsible global citizen---is high.

But Germany also burns lignite, the dirtiest coal. It imports Russian gas. \
It shut down nuclear plants while coal plants stayed open. When asked, \
there are good reasons: energy security, jobs, political coalitions.

Germany optimizes its national coherence, $K_{comm}$. It feels like a \
climate leader.

\paragraph{Level 3: The System.}

The atmosphere doesn't care about feelings.

Global CO\textsubscript{2} continues to accumulate. Temperatures rise. Ice \
melts. Seas rise. Storms intensify. The climate system moves toward states \
not seen in human history.

System coherence, $K_{sys}$, declines even as individual and national \
coherence increase.

\subsubsection{The Mathematical Structure}

The framework captures this with precision:

\begin{equation}
K_{ind}^* \neq K_{comm}^* \neq K_{sys}^*
\end{equation}

When externalities cross levels, optimizing at any single level doesn't \
optimize the whole. The Nash equilibrium---where each level optimizes its \
own coherence---produces:

\begin{equation}
K_{total} < K_{cooperative}
\end{equation}

This isn't a bug. It's the structure.

\subsubsection{Why Awareness Makes It Worse}

Here's the twist: higher awareness can deepen the trap.

When Maria becomes more aware of climate change ($A_M^{climate} \uparrow$), \
she feels more anxious. To restore coherence, she takes visible personal \
actions---reusable bags, electric car, carbon offsets.

These actions increase $K_{ind}$ without changing $K_{sys}$. Worse, they \
can reduce political pressure for systemic change, because people feel \
like they're ``doing something.''

The same pattern operates at the national level. Germany's renewable \
energy investments create coherence (``we're climate leaders'') that \
reduces pressure for harder actions (ending coal, rethinking industrial \
policy).

Awareness → Individual action → Personal coherence → Reduced systemic pressure

This is AWX-A1a and AWX-A1b in tension: self-referential awareness produces \
one response, community-referential awareness would produce another.

\subsubsection{The Numbers}

From the Γ-matrix estimation:

\begin{center}
\renewcommand{\arraystretch}{1.3}
\begin{tabular}{lcc}
\hline
\textbf{Parameter} & \textbf{Value} & \textbf{Interpretation} \\
\hline
$\gamma_{S,S}$ & +0.73 & Social context amplifies social utility \\
$w_{IDN}/w_{INU}$ & 2.1 & Identity weighs more than material outcomes \\
$\lambda_{cross-level}$ & 0.24 & Only 24\% of effects transmit up \\
\hline
\end{tabular}
\end{center}

The cross-level transmission coefficient (0.24) is the key number. It \
means that individual actions have only a fractional effect on system \
outcomes---but they have a full effect on individual coherence.

Rational agents therefore overinvest in individual coherence relative to \
system outcomes.

\subsubsection{Testable Implication}

\begin{tcolorbox}[colback=blue!5!white, colframe=blue!75!black, title=Prediction 7: Multi-Level Conflict]
Climate policy outcomes will show systematic divergence between individual, \
national, and global optimality, with each level appearing ``rational'' \
at its own scale while producing suboptimal aggregate outcomes.
\end{tcolorbox}

Specific test: Countries with highest individual climate awareness will \
not show proportionally higher emission reductions, because awareness \
drives personal coherence-restoration rather than systemic change.

\subsubsection{What This Means}

The multi-level coherence conflict predicts that:

\begin{enumerate}
    \item Individual virtue alone cannot solve collective problems
    \item National leadership alone cannot solve global problems
    \item Coherence at one level can undermine coherence at another
    \item Effective policy must align incentives \emph{across} levels
\end{enumerate}

The framework suggests that successful climate policy requires explicit \
attention to coherence transmission: mechanisms that make system-level \
outcomes salient for individual-level coherence.

Carbon pricing does this, imperfectly. It makes individual choices \
contribute to system outcomes in visible, traceable ways. But it works \
only if the signal is strong enough to penetrate individual awareness \
and reshape coherence calculations.

The Paris Agreement failed this test. It created national coherence \
(``we signed!'') without individual salience (``what does it mean for me?'') \
or system accountability (``who checks?'').

\subsubsection{The Deeper Insight}

Elinor Ostrom won the Nobel Prize for showing that commons problems can \
be solved---but only with careful institutional design that aligns \
incentives across levels.

Our framework formalizes why: the coherence function must be specified \
at the system level, not imported from individual or community levels.

This is prediction 7: when everyone optimizes their own coherence, no \
one gets what they want.

\begin{cross-reference}
\textbf{Formal Specification:} PRED-07 with mathematical treatment in \
Appendix AWA, Section~\ref{subsec:pred-07-calc} \
(Novel Predictions: Formal Specification).
\end{cross-reference}
